% RESULTADO_RECUPERADO: c40_termicas.tex del integral activo.
% Ampliación de prueba: conteo de modos, momentos térmicos, máximos
% espectrales y sustitución de la escala del calendario.
% Las hipótesis de realización se conservan explícitas.
\section{Radiación térmica y escala de información del calendario}
\label{v:sec:radiacion-termica}

La estadística de ocupación proporciona el enlace entre los modos de
una realización y su energía media. En la realización electromagnética,
la dispersión y la multiplicidad de polarización determinan además la
densidad espectral. Su composición permite obtener los momentos térmicos
y los máximos espectrales, manteniendo separadas tres operaciones:
la producción de la escala de acción, la representación de los modos y
la evaluación del estado térmico.

\subsection{Acción, calendario y temperatura}

Se conserva la orientación \(a\) del reloj y se abrevia
\(\hbar=\hbar_a\), \(c=c_{\rm int}=c_{\rm vac}\), con la
identificación de \eqref{v:eq:identidad-velocidad} en una misma coordenatización.
La energía \(\varepsilon_{108}\) de esta sección es
\(\epsilon_a\) de \eqref{v:ant:eq:ii-kb-aut-energia};
\(\Theta_{108}=\Theta_{\rm clk,a}\). Su división por \(\log3\)
produce la energía por nat de \eqref{v:ant:eq:ii-kb-aut-energia-nat}.
La construcción y la covariancia de la pareja térmica se encuentran en
la sección~\ref{v:ant:sec:ii-kb-aut-desarrollo}.

Escribimos \(h=2\pi\hbar>0\) para la representación de acción, \(t_0>0\)
para la unidad temporal de su coordenatización dimensional y \(\ell_0=ct_0\) para
la longitud correspondiente. La celda cronológica de longitud \(108\)
determina la energía
\begin{equation}
 \varepsilon_{108}=\frac{h}{108t_0}
 =\frac{\hbar c}{L_\alpha}=Q_{L,a}=E_{P,a},\qquad
 k_B\Theta_{108}\log 3=Q_{L,a},\qquad
 \beta_{108}=\frac{108t_0\log3}{h}.
 \label{v:eq:energia-informacion-calendario}
\end{equation}
Aquí \(\beta=(k_B\Theta)^{-1}\) es el parámetro de Gibbs. La igualdad
fija el producto energético de la temperatura y el transductor de
entropía. La individualización dimensional de \(k_B\) y
\(\Theta_{108}\) utiliza la sección térmica y sus bases. El logaritmo
natural mide información por rama trítica; su factor \(\log3\) convierte
la cardinalidad de la celda en entropía adimensional.

Introducimos también las unidades asociadas
\begin{equation}
 t_P=\frac{108}{2\pi_{\HMT}}t_0=\frac{L_\alpha}{c},\qquad
 \ell_P=ct_P=L_\alpha,\qquad
 E_P:=E_{P,a}=\frac{\hbar}{t_P}=Q_{L,a}=\varepsilon_{108}.
 \label{v:eq:unidades-termicas-calendario}
\end{equation}
Estas igualdades precisan el significado de los símbolos en esta sección;
su evaluación en otra coordenatización conserva las identidades entre cocientes.
La longitud \(\ell_P=(54/\pi_{\HMT})\ell_0\) coincide primero con
la longitud estructural \(L_\alpha\) de
\eqref{v:eq:longitud-alpha}; la identidad gravitatoria posterior
\(\ell_P^2=G_a\hbar_a/c^3\) es el corolario
\eqref{v:eq:gravedad-unica-alpha}, no el origen de esa escala.

La equivalencia centrada
\eqref{v:eq:equivalencia-termica-centrada} y la identidad del transductor
\eqref{v:eq:identidad-transductor-numeral} fijan, antes de la realización
electromagnética, la escala energética y su lectura térmica. El marcador
$p_0$ de \eqref{v:eq:particion-portador-termico} conserva la normalización
y el origen al condicionar una rama. Esta sustitución determina $\beta$ y
la escala de energía; no selecciona retrospectivamente la dispersión ni
las dos polarizaciones que se especifican a continuación. Por ello las
leyes de Planck, Stefan--Boltzmann y Wien permanecen como consecuencias
posteriores de la realización modal declarada.

\subsection{Densidad de modos y ley espectral}

\begin{proposition}[Realización térmica de modos transversales]
\label{v:prop:ley-espectral}
Considérese una realización libre en una caja periódica tridimensional,
cuyo espacio de excitaciones contiene los vectores de onda
\(\mathbf k\ne0\), con dos modos transversales por vector, dispersión
\(\omega=c|\mathbf k|\), energía de excitación \(\hbar\omega\) y
potencial químico nulo y \(\beta>0\). El modo constante se mantiene
fijado como dato de fondo. En el límite termodinámico, la densidad de energía
por frecuencia angular es
\begin{equation}
 u_\omega(\Theta)=\frac{\hbar\omega^3}{\pi^2c^3}
 \frac{1}{e^{\beta\hbar\omega}-1},\qquad \omega>0.
 \label{v:eq:planck-angular}
\end{equation}
La energía es la de las excitaciones sobre el vacío de esta
representación.
\end{proposition}

\begin{proof}
Para una caja de volumen \(V\), los vectores de onda tienen densidad
\(V/(2\pi)^3\). En una corona radial de grosor \(dk\), la densidad
por volumen, incluyendo las dos polarizaciones, es
\[
 2\frac{4\pi k^2\,dk}{(2\pi)^3}=\frac{k^2}{\pi^2}\,dk
 =\frac{\omega^2}{\pi^2c^3}\,d\omega.
\]
Se trata de la densidad límite de las sumas sobre la red periódica
sin el modo constante. Esta elección define el producto de partición
antes del límite: cada modo de excitación tiene energía positiva.
La exclusión de un único punto conserva la densidad límite radial.
La ocupación bosónica previamente obtenida es
\(\bar n_\omega=(e^{\beta\hbar\omega}-1)^{-1}\).
Multiplicar la densidad por \(\hbar\omega\bar n_\omega\) demuestra
\eqref{v:eq:planck-angular}. Cerca de cero el integrando energético es
de orden \(\omega^2\), y en infinito decrece exponencialmente, de modo
que su integral converge.
\end{proof}

La densidad procede de la representación tridimensional especificada.
La multiplicidad transversal utilizada aquí tiene ese dominio físico;
la dimensión de un espacio armónico sobre el toro APP describe, por
separado, la cohomología de su complejo celular.

% FORMALIZACION_NUEVA de la justificación del límite utilizado en
% 70_radiacion_termica.tex. No introduce un generador de constantes.
\subsection{Límite termodinámico de las sumas de ocupación}
\label{v:subsec:limite-termodinamico}

La densidad radial anterior puede obtenerse como límite efectivo de las
sumas de la realización periódica. La estimación debe conservar el modo
constante fijado y controlar uniformemente tanto el origen como la cola
espectral. Las representaciones de acción y calendario permanecen previas
a esta operación de realización.

\begin{lemma}[Convergencia de los tres observables térmicos]
\label{v:lem:limite-termodinamico}
Sean \(a,b>0\), \(L>0\), \(\Delta=2\pi/L\) y
\[
 F_N(r)=\frac{1}{e^{ar}-1},\qquad
 F_E(r)=\frac{br}{e^{ar}-1},\qquad
 F_Z(r)=-\log(1-e^{-ar})\quad(r>0).
\]
Para cada una de estas funciones se cumple
\begin{equation}
 \lim_{L\to\infty}\frac{2}{L^3}
 \sum_{\nu\in\mathbb Z^3\setminus\{0\}}
 F\!\left(\frac{2\pi}{L}|\nu|\right)
 =\frac{1}{\pi^2}\int_0^\infty r^2F(r)\,dr.
 \label{v:eq:limite-modos-discretos}
\end{equation}
Las sumas y la integral son finitas. La especialización
\(a=\beta\hbar c\), \(b=\hbar c\) proporciona, respectivamente,
las densidades de número, energía de excitación y logaritmo de la función
de partición de los modos transversales.
\end{lemma}

\begin{proof}
Las tres funciones son continuas y positivas en \((0,\infty)\).
Existen constantes \(C_0,C_1,d>0\), dependientes de \(a,b\) pero
independientes de \(\Delta\), tales que
\begin{equation}
 F(r)\leq\frac{C_0}{r}\quad(0<r\leq1),\qquad
 F(r)\leq C_1e^{-dr}\quad(r\geq1).
 \label{v:eq:envolventes-termicas}
\end{equation}
En efecto, \(e^{ar}-1\geq ar\) controla \(F_N\) y \(F_E\)
cerca del origen; además,
\(F_Z(r)=\log(1+F_N(r))\leq F_N(r)\leq1/(ar)\).
Para \(r\geq1\), la desigualdad
\(-\log(1-y)\leq y/(1-y)\), con \(y=e^{-ar}\), da la cota
exponencial de \(F_Z\). Las otras dos siguen de
\((e^{ar}-1)^{-1}\leq e^{-ar}/(1-e^{-a})\), absorbiendo el factor
\(r\) de \(F_E\) en una exponencial de menor tasa. Estas cotas
prueban ya la integrabilidad de \(r^2F(r)\).

Para las sumas discretas, agrupamos los puntos enteros por
\(m=\|\nu\|_\infty\geq1\). Cada capa contiene exactamente
\[
 (2m+1)^3-(2m-1)^3=24m^2+2
\]
puntos y satisface \(m\leq|\nu|\leq\sqrt3m\).
Fijemos \(0<\delta\leq1\) y \(0<\Delta\leq1\).
Si \(\delta<\Delta\), no hay puntos con
\(0<\Delta|\nu|\leq\delta\). En otro caso, poniendo
\(N=\lfloor\delta/\Delta\rfloor\), la primera cota de
\eqref{v:eq:envolventes-termicas} permite escribir
\begin{align*}
 \Delta^3\!\sum_{0<\Delta|\nu|\leq\delta}F(\Delta|\nu|)
 &\leq C_0\Delta^2\sum_{m=1}^{N}\frac{24m^2+2}{m}\\
 &\leq26C_0\Delta^2\sum_{m=1}^{N}m
 \leq26C_0\delta^2.
\end{align*}
En la segunda línea se ha usado \(m^{-1}\leq m\) y, después,
\(N(N+1)/2\leq N^2\) para \(N\geq1\).
La contribución integral del mismo entorno también es de orden
\(\delta^2\), pues
\(\int_0^\delta r^2F(r)\,dr\leq C_0\delta^2/2\).

Para \(R\geq1\), todo punto con \(\Delta|\nu|>R\) pertenece
a una capa \(m>R/(\sqrt3\Delta)\). La segunda cota de
\eqref{v:eq:envolventes-termicas} da
\begin{align*}
 \Delta^3\!\sum_{\Delta|\nu|>R}F(\Delta|\nu|)
 &\leq C_1e^{-dR/(2\sqrt3)}
 \Delta^3\sum_{m\geq1}(24m^2+2)e^{-d\Delta m/2}\\
 &\leq C_2e^{-dR/(2\sqrt3)},
\end{align*}
donde \(C_2\) no depende de \(0<\Delta\leq1\). Para comprobar
esa uniformidad, se toma \(q=e^{-d\Delta/2}\) y se utilizan
\[
 \sum_{m\geq1}q^m=\frac{q}{1-q},\qquad
 \sum_{m\geq1}m^2q^m=\frac{q(1+q)}{(1-q)^3}.
\]
La primera identidad procede de la serie geométrica y la segunda de
aplicar dos veces \(q\,d/dq\), operación válida en \(|q|<1\).
Además,
\(1-e^{-d\Delta/2}\geq\Delta(1-e^{-d/2})\) por concavidad.
El factor \(\Delta^3\) compensa, por tanto, los denominadores de
las dos expresiones. La cola integral tiende asimismo a cero por la
cota exponencial. Para cada \(\Delta>0\), estas mismas estimaciones
prueban la convergencia de la suma infinita.

En el anillo compacto \(\delta\leq|k|\leq R\), la función
\(F(|k|)\) es continua y acotada. Las sumas de Riemann de malla
\(\Delta\) convergen a su integral; las dos esferas de frontera
tienen medida nula. Se hace primero \(\Delta\to0\), después
\(\delta\to0\) y finalmente \(R\to\infty\), usando las cotas
uniformes anteriores. Resulta
\[
 \Delta^3\sum_{\nu\ne0}F(\Delta|\nu|)
 \longrightarrow \int_{\mathbb R^3}F(|k|)\,d^3k
 =4\pi\int_0^\infty r^2F(r)\,dr.
\]
Como \(2/L^3=2\Delta^3/(2\pi)^3\), el coeficiente radial final
es \(2\cdot4\pi/(2\pi)^3=1/\pi^2\), lo que demuestra
\eqref{v:eq:limite-modos-discretos}.
\end{proof}

El control del origen se efectúa sobre excitaciones de energía positiva.
El modo constante permanece fijado antes de formar la función de partición:
su eliminación posterior de una integral no sustituiría esa condición de
la representación finita. La prueba establece el paso al límite de la
realización declarada y conserva la distinción entre esta realización y
la construcción anterior de sus parámetros de acción y calendario.


\subsection{Momentos térmicos e información}

\begin{lemma}[Momento bosónico]
\label{v:lem:momento-bosonico}
Para todo real \(s>1\),
\begin{equation}
 \int_0^\infty\frac{x^{s-1}}{e^x-1}\,dx
 =\Gamma(s)\zeta(s).
 \label{v:eq:momento-bosonico}
\end{equation}
\end{lemma}
\begin{proof}
Para \(x>0\), \((e^x-1)^{-1}=\sum_{m\geq1}e^{-mx}\).
Todos los sumandos son no negativos; Tonelli permite intercambiar suma
e integral. El cambio \(y=mx\) da
\(\int_0^\infty x^{s-1}e^{-mx}\,dx=m^{-s}\Gamma(s)\).
La suma converge para \(s>1\), y resulta la igualdad anunciada.
\end{proof}

\begin{proposition}[Densidades de número, energía y entropía]
\label{v:prop:densidades-termicas}
En la realización de la proposición~\ref{v:prop:ley-espectral},
\begin{align}
 n_\gamma&=\frac{2\zeta(3)}{\pi^2c^3(\beta\hbar)^3},\\
 u_\gamma&=\frac{\pi^2}{15c^3\hbar^3\beta^4},\\
 s_\gamma&=\frac{4\pi^2 k_B}{45c^3(\beta\hbar)^3}.
 \label{v:eq:densidades-termicas}
\end{align}
En la temperatura del calendario, estas relaciones se escriben
\begin{equation}
 n_\gamma\ell_P^3=\frac{2\zeta(3)}{\pi^2(\log3)^3},\quad
 \frac{u_\gamma\ell_P^3}{E_P}=\frac{\pi^2}{15(\log3)^4},\quad
 \frac{s_\gamma\ell_P^3}{k_B}=\frac{4\pi^2}{45(\log3)^3}.
 \label{v:eq:densidades-calendario}
\end{equation}
\end{proposition}

\begin{proof}
La integral de la ocupación por la densidad de modos usa
\eqref{v:eq:momento-bosonico} con \(s=3\), y la integral energética
lo usa con \(s=4\). Se obtiene \(\Gamma(3)=2\) y
\(\Gamma(4)\zeta(4)=6\pi^4/90=\pi^4/15\).
Para la entropía se parte del producto de partición de los modos:
\[
 \frac{\log Z}{V}=-\frac{1}{\pi^2c^3}
 \int_0^\infty\omega^2
 \log(1-e^{-\beta\hbar\omega})\,d\omega
 =\frac{\pi^2}{45c^3(\beta\hbar)^3}.
\]
La última igualdad se obtiene integrando por partes tras el cambio
\(x=\beta\hbar\omega\). Los términos de frontera se anulan:
\(x^3\log(1-e^{-x})\) tiende a cero en ambos extremos.
La identidad de Gibbs \(s_\gamma/k_B=\log Z/V+\beta u_\gamma\)
prueba la tercera expresión. Finalmente,
\(\beta_{108}\hbar=t_P\log3\) y \(\ell_P=ct_P\) dan las tres
normalizaciones del calendario.
\end{proof}

\subsection{Composición fotónica del nivel \texorpdfstring{\(A_2\)}{A2} de Bendersky--Glaisher}
\label{v:subsec:bendersky-fotonica}

La densidad de número de la proposición~\ref{v:prop:densidades-termicas}
admite una segunda lectura exacta que conserva su procedencia espectral.
Sea \(\mathcal Z_3\) el funcional escalar normalizado producido por el
lector espectral trítico y reconocido posteriormente como la función zeta
en la realización arquimediana. Sean
\[
 H_k=\sum_{j=1}^{k}\frac1j,\qquad H_0=0,
\]
y \(\mathsf B_j\) los números de Bernoulli con el convenio
\(\mathsf B_1=-1/2\). Para evitar toda colisión con la notación de redes,
definimos los niveles de Bendersky--Glaisher mediante
\begin{equation}
 A_k^{\mathrm{BG}}
 =\exp\!\left(
   \frac{H_k\mathsf B_{k+1}}{k+1}-\mathcal Z_3'(-k)
 \right),\qquad k\in\mathbb N_0.
 \label{v:eq:bendersky-bg}
\end{equation}
El término de Bernoulli fija la normalización del producto regularizado;
no introduce una constante física. La coordenada
\(\pi_{\mathrm{HMT}}\) que interviene en la ecuación funcional ya fue
expresada por el estado enriquecido APP--TRIT--TPK antes de esta lectura.

\begin{proposition}[Composición fotónica del nivel dos]
\label{v:prop:bendersky-fotonica}
En la normalización escalar \(\mathcal Z_3=\zeta\),
\begin{equation}
 \log A_2^{\mathrm{BG}}
 =-\mathcal Z_3'(-2)
 =\frac{\mathcal Z_3(3)}{4\pi_{\mathrm{HMT}}^2}.
 \label{v:eq:bendersky-nivel-dos}
\end{equation}
Por consiguiente, en la coordenatización térmica del calendario,
\begin{align}
 n_\gamma\ell_P^3
  &=\frac{8\log A_2^{\mathrm{BG}}}{(\log3)^3}
   =-\frac{8\mathcal Z_3'(-2)}{(\log3)^3},
   \label{v:eq:numero-fotones-bendersky}\\
 \frac{u_\gamma}{n_\gamma k_B\Theta_{108}}
  &=\frac{\pi_{\mathrm{HMT}}^2}{120\log A_2^{\mathrm{BG}}},
   \label{v:eq:energia-media-bendersky}\\
 \frac{s_\gamma}{n_\gamma k_B}
  &=\frac{\pi_{\mathrm{HMT}}^2}{90\log A_2^{\mathrm{BG}}}.
   \label{v:eq:entropia-media-bendersky}
\end{align}
\end{proposition}

\begin{proof}
La forma completada satisface
\[
 \pi_{\mathrm{HMT}}^{-s/2}\Gamma(s/2)\mathcal Z_3(s)
 =\pi_{\mathrm{HMT}}^{-(1-s)/2}
  \Gamma((1-s)/2)\mathcal Z_3(1-s).
\]
En \(s=-2\), el cero simple de \(\mathcal Z_3\) cancela el polo de
\(\Gamma(s/2)\). Comparar los términos constantes, usando
\(\Gamma(3/2)=\sqrt{\pi_{\mathrm{HMT}}}/2\) tras el reconocimiento
arquimediano, da
\[
 \mathcal Z_3'(-2)=-\frac{\mathcal Z_3(3)}{4\pi_{\mathrm{HMT}}^2}.
\]
Como \(\mathsf B_3=0\), la definición~\eqref{v:eq:bendersky-bg}
produce la primera igualdad de~\eqref{v:eq:bendersky-nivel-dos}.
Sustituir
\(\mathcal Z_3(3)=4\pi_{\mathrm{HMT}}^2\log A_2^{\mathrm{BG}}\)
en~\eqref{v:eq:densidades-calendario} demuestra
\eqref{v:eq:numero-fotones-bendersky}. Finalmente,
\(k_B\Theta_{108}=E_P/\log3\); dividir las densidades de energía y
entropía por la densidad de número da
\eqref{v:eq:energia-media-bendersky} y
\eqref{v:eq:entropia-media-bendersky}.
\end{proof}

La evaluación \(A_2^{\mathrm{BG}}=1.0309167521973921\ldots\) sirve como
test numérico de las tres identidades, no como dato generador. El desarrollo
espectral Bendersky--Glaisher contiene la jerarquía completa de derivadas
regularizadas y su correspondencia con los valores espectrales impares; la presente sección
establece la composición recíproca que utiliza específicamente la densidad
adimensional de número fotónico. No se redefine \(k_B\), no se identifica una densidad con un
operador regularizado y no se altera ninguna de las densidades con
\(\zeta(3)\) ya demostradas.


\subsection{Flujo térmico y constante de Stefan--Boltzmann}

\begin{proposition}[Flujo isotrópico y normalización de acción]
\label{v:prop:stefan}
El flujo hemisférico de la realización isotrópica es
\begin{equation}
 j_*(\Theta)=\sigma\Theta^4,\qquad
 \sigma=\frac{\pi^2 k_B^4}{60\hbar^3c^2}
 =\frac{2\pi^5k_B^4}{15h^3c^2}.
 \label{v:eq:stefan}
\end{equation}
En la temperatura del calendario,
\begin{equation}
 j_*(\Theta_{108})
 =\frac{2\pi^5}{15\,108^4(\log3)^4}
 \frac{h}{\ell_0^2t_0^2}.
 \label{v:eq:stefan-calendario}
\end{equation}
\end{proposition}
\begin{proof}
La energía isotrópica por unidad de ángulo sólido es
\(u_\gamma/(4\pi)\). El flujo normal a una superficie integra
\(c\cos\theta\) sobre el hemisferio:
\[
 \frac{cu_\gamma}{4\pi}
 \int_0^{2\pi}\!d\phi\int_0^{\pi/2}\cos\theta\sin\theta\,d\theta
 =\frac{cu_\gamma}{4}.
\]
La fórmula para \(u_\gamma\), \(\beta^{-1}=k_B\Theta\) y
\(h=2\pi\hbar\) demuestran \eqref{v:eq:stefan}.
Sustituir \eqref{v:eq:energia-informacion-calendario} y
\(c=\ell_0/t_0\) produce \eqref{v:eq:stefan-calendario}.
\end{proof}

La constante \(\sigma\) es, dentro de esta realización, el coeficiente
que compone la ocupación bosónica, la densidad de modos transversales y
la proyección angular del flujo. Su potencia cuarta procede del momento
cúbico de frecuencia junto con el diferencial de integración. La lectura
del calendario expresa ese mismo flujo mediante acción, longitud,
tiempo e información trítica.

\newpage
\subsection{Máximos espectrales y ley de desplazamiento}

\begin{lemma}[Raíz positiva de los máximos espectrales]
\label{v:lem:raiz-wien}
Para todo entero \(n>1\), la ecuación
\(n(1-e^{-x})=x\) tiene una única solución \(x_n>0\).
La función \(f_n(x)=x^n/(e^x-1)\) tiene su máximo positivo único en
\(x_n\).
\end{lemma}
\begin{proof}
Sea \(g_n(x)=n(1-e^{-x})-x\). Se cumple
\(g_n(0)=0\), \(g_n'(x)=ne^{-x}-1\) y
\(g_n''(x)=-ne^{-x}<0\). El máximo de \(g_n\) se alcanza en
\(\log n\), donde su valor es \(n-1-\log n>0\), mientras
\(g_n(n)=-ne^{-n}<0\). La monotonía estricta a cada lado del máximo
prueba existencia y unicidad de la raíz positiva. El signo de
\(f_n'(x)\) coincide con el de \(g_n(x)\); además, \(f_n\) tiende
a cero en cero y en infinito. Así queda probado el máximo único.
\end{proof}

\begin{proposition}[Lecturas en frecuencia y longitud de onda]
\label{v:prop:wien}
Los máximos de radiancia por frecuencia y por longitud de onda son
\begin{equation}
 \nu_{\max}=\frac{x_3 k_B\Theta}{h},\qquad
 \lambda_{\max}=\frac{hc}{x_5 k_B\Theta}.
 \label{v:eq:wien}
\end{equation}
En particular,
\begin{equation}
 \nu_{\max}(\Theta_{108})=
 \frac{x_3}{108t_0\log3},\qquad
 \lambda_{\max}(\Theta_{108})=
 \frac{108\log3}{x_5}\ell_0.
 \label{v:eq:wien-calendario}
\end{equation}
\end{proposition}
\begin{proof}
La conversión \(\omega=2\pi\nu\) y el factor angular isotrópico dan
\[
 B_\nu=\frac{2h\nu^3}{c^2}\frac1{e^{\beta h\nu}-1}.
\]
Con \(\nu=c/\lambda\), la densidad transformada satisface
\(B_\lambda=B_\nu|d\nu/d\lambda|\); por tanto,
\[
 B_\lambda=\frac{2hc^2}{\lambda^5}
 \frac1{e^{\beta hc/\lambda}-1}.
\]
Los cambios de variable respectivos conducen a \(f_3\) y \(f_5\).
El lema anterior prueba \eqref{v:eq:wien}; la escala del calendario
da \eqref{v:eq:wien-calendario}. El jacobiano de densidad explica que
los dos máximos correspondan a raíces distintas:
\(\nu_{\max}\lambda_{\max}=c x_3/x_5\).
\end{proof}

El desplazamiento de Wien queda definido estructuralmente por el punto
estacionario de la densidad espectral elegida. El cambio de variable
transforma tanto el argumento como la medida. Esta precisión sitúa
cada constante espectral en su operación generadora y conserva su
significado físico al pasar entre frecuencia y longitud de onda.
